Learning-based methods augment SBMP by learning a context-conditioned distribution, thereby concentrating samples in feasible and task-relevant regions of the state space.
Early probabilistic learning approaches have been used to bias SBMP via Gaussian Mixture Models (GMMs) learned online from collision and collision-free exemplars, enabling sample biasing toward predicted free regions~\cite{Jinwook2016GMM, Jinwook2017GMM}. However, the method depends on continual acquisition of informative exemplars and careful tuning of parameters.
In contrast, the CVAE is trained offline on a large set of traffic scenarios and requires only a single forward pass conditioned on the current BEV at runtime, without online collection of exemplars or online adaptation of the distribution parameters.
\textcite{Ma2020ConditionalGA} employ a conditional Generative Adversarial Network (GAN) that operates on rasterized scene representations and predicts dense likelihood maps.
These likelihood maps require additional post-processing to extract planner states and remain tightly coupled with geometric planners, as temporal dynamics are not explicitly modeled.
The CVAE, by contrast, directly outputs the Fren\'{e}t terminal parameters $(d, \dot{s}, T)$, which are passed to the polynomial trajectory generation of FOP without post-processing, while temporal information is provided by the three stacked BEV frames and the planning horizon $T$. Moreover, the CVAE is trained by maximizing the ELBO and does not require adversarial training.
Transformer-based sampling methods~\cite{Johnson2023Dictionaries, Lei2024TransformerEnhancedMP, Feng2025rrtformer} have also been integrated into SBMP by conditioning sampling proposals on the planning context, but they are typically invoked sequentially at inference time, which scales poorly when large batches of independent samples are required.
The CVAE instead encodes the conditioning images once and generates all $N$ samples in parallel within a single forward pass by drawing a batch of latent variables $z$ from the prior, so that its inference time grows only marginally with $N$.
Related ideas appear in FOP, where terminal states in the Fren\'{e}t frame are sampled, and then the trajectories are generated by polynomials. LF-Net~\cite{Yu2024LFNet} follows this structure by uniformly discretizing potential terminal states and scoring them with a cross-attention-based classifier. This enables selecting a small set of high-probability candidates before running full trajectory optimization.
Unlike per-sample trajectory evaluation, this reduces computation by shifting most of the candidate scoring process to a parallelizable scoring stage, but it remains tied to discretization since solution quality depends on the range and resolution of the terminal-state lattice.
The CVAE, in contrast, generates the terminal parameters in a continuous parameter space that is not limited by a lattice resolution, and the number of samples $N$ can be chosen freely according to the computational budget.

\textcite{Ichter2018SamplingDistributions} propose a non-uniform sampling strategy for SBMP that favors sampling in those regions where an optimal solution is more likely to lie. A CVAE is used to learn this sampling distribution, conditioned on information specific to each planning problem. The method demonstrates flexibility and effectiveness by learning optimal, conditioned sampling distributions from demonstrations for robotics tasks such as static obstacle avoidance, path planning, and overtaking maneuvers.
The approach presented in this paper builds on this line of work but generalizes it to the domain of autonomous driving. In particular, the BEV-based conditioning allows to handle arbitrary traffic scenarios, rather than being limited to specific maneuvers or predefined subtasks.
Furthermore, the CVAE replaces only the sampling step of FOP: the coordinate transformation, the kinematic constraint checks, the collision check, and the cost evaluation remain unchanged, and FOP with uniform sampling serves as a fallback, so that safety does not depend on the network output.

Diffusion models have recently also been applied to planning for autonomous driving and robotics~\cite{Yang2024DiffusionES, Liao2025DiffusionDrive}. Their iterative denoising process provides high expressiveness, but each generation requires multiple network evaluations.
For example, even with a reduced number of denoising steps and a smaller population, Diffusion-ES still requires about 500~ms per planning cycle, which far exceeds the planning cycle of 0.1~s considered in this work.
Subsequent work reduces the number of steps by high-order solvers or truncated denoising~\cite{Liao2025DiffusionDrive}, but still relies on dedicated acceleration techniques.
In contrast, the CVAE generates all candidate samples with a single forward pass.

% --- BibTeX entries for the two new citations (verified against the arXiv PDFs) ---
% @inproceedings{Yang2024DiffusionES,
%   author    = {Yang, Brian and Su, Huangyuan and Gkanatsios, Nikolaos and Ke, Tsung-Wei and Jain, Ayush and Schneider, Jeff and Fragkiadaki, Katerina},
%   title     = {Diffusion-{ES}: Gradient-free Planning with Diffusion for Autonomous and Instruction-guided Driving},
%   booktitle = {Proceedings of the IEEE/CVF Conference on Computer Vision and Pattern Recognition (CVPR)},
%   year      = {2024}
% }
% @inproceedings{Liao2025DiffusionDrive,
%   author    = {Liao, Bencheng and Chen, Shaoyu and Yin, Haoran and Jiang, Bo and Wang, Cheng and Yan, Sixu and Zhang, Xinbang and Li, Xiangyu and Zhang, Ying and Zhang, Qian and Wang, Xinggang},
%   title     = {{DiffusionDrive}: Truncated Diffusion Model for End-to-End Autonomous Driving},
%   booktitle = {Proceedings of the IEEE/CVF Conference on Computer Vision and Pattern Recognition (CVPR)},
%   year      = {2025}
% }
